\section{Discussion}
\label{sec:discussion}
\para{What does isolation mean here?} The best tested primary preservation condition still changes 4.05\% of protected answers. Broader training expands semantic transfer while reaching additional protected behavior. These measurements support a scope-dependent trade-off for the tested updates; they do not establish that all faithful semantic edits must damage unrelated behavior. The wrapper supplies a different result: outside its explicit key, the base execution is unchanged. That architectural guarantee neither updates arithmetic computation nor demonstrates implanted belief.

\para{What limits the inference?} We train one small historical model at one prospectively fixed layer. We do not sweep layers, run full-model fine-tuning, causally patch arithmetic circuits, or evaluate specialized editors. The auxiliary API study does not remedy the absence of a second weight-access model. Three seeds provide coarse inference, and plain \ft repeats are deterministic. Secondary identities are case checks rather than a random sample of edits.

The finite suite contains only 17 main semantic surfaces and 60 dependent expressions from three operator families. Cluster bootstrapping partly addresses shared templates but cannot establish universal preservation. Baseline performance also varies sharply: \cfact has 26/128 strict matches and \gsm only 1/32. Incomplete factual prompts, aliases, and a 16-token answer-only budget constrain those scores. These sets primarily measure output stability here, not robust retention of general capability.

Answer/end-of-turn supervision changes formatting as well as answers. The post-hoc matched sham exposes this confound. Breadth matches updates and supervised tokens but uses roughly 9--10\% more edit input tokens. Replay covers a narrow operand range and contains no \wikitext passages. Other replay distributions, regularization strengths, layers, or checkpoints might improve preservation. We evaluate the earliest successful checkpoint rather than a prospectively optimized efficacy--locality boundary.

\para{What can adoption tell us about belief?} Our local rule avoids assigning contradictory global arithmetic labels, but models can use strategies that bypass the edited subexpression. Some original errors match counterfactual answers. Weak dependent transfer and strong interface sensitivity therefore matter more than exact-key success alone. We cannot infer latent belief, causal representations, or stable acceptance across contexts. The parser, zero-operand analysis, format audit, and 20-step sham are explicitly post-hoc; the two predeclared contrasts remain unchanged.

\para{Broader implications.} Editors used for factual maintenance or unlearning should specify their intended propagation scope and report output changes separately from errors. A targeted false response can also serve deceptive uses; our deliberately false arithmetic task is a diagnostic intervention, not a recommendation to deploy misinformation. Explicit gates offer inspectable scope, whereas semantic updating requires evaluation across wording, consequences, and interfaces. Neither a finite zero-change score nor a small language-model KL establishes that a deployed system is generally unaffected.
